% Consolidated guide-ch08 from 9a1384ddf1afcd855b09256c9c6f1c90b5023140; source fragments preserved.

% === BEGIN guide-ch08-examples.tex blob 942b76434e360dcd15f0f8e0eaaf17a6c8f86b03 ===
% 研究性题解：赵近芳、王登龙《大学物理学学习指导》第5版，第8章“热力学基础”
% 例8.1—8.6，原书印刷111—114页 / 扫描PDF物理119—122页，2026-11-23逐页查原扫描。
% 题意以短摘要转述，所有推导独立编写；本批待独立学术二校。

\sourceexample{8.1}{111--112}{119--120}
\problemtext 标准状态下$0.014\,\mathrm{kg}$的氮气（刚性双原子理想气体）
经(1)等温、(2)等压两种过程压缩至原体积的一半，求各过程内能变化、传热和外界对气体做功。
\workings
摩尔量$n=0.014/0.028=0.5\,\mathrm{mol}$，标准温度$T_1=273\,\mathrm K$。
用气体\emph{对外做功}$W$为正的约定：$\Delta U=Q-W$。
(1) 等温时$\Delta U_T=0$且
\[
W_T=nRT_1\ln\frac{V_2}{V_1}
=\frac12 R(273)\ln\frac12\approx-7.86\times10^2\,\mathrm J.
\]
因此
\[
\boxed{\Delta U_T=0},\quad
\boxed{Q_T\approx-786\,\mathrm J},\quad
\boxed{W_{\mathrm{on},T}=-W_T\approx786\,\mathrm J}.
\]
(2) 等压时由理想气体方程$T_2/T_1=V_2/V_1=1/2$，
$\Delta T=-136.5$ K，故
\[
\boxed{\Delta U_p=n\frac52 R\Delta T\approx-1.418\times10^3\,\mathrm J},
\quad
\boxed{W_p=nR\Delta T\approx-567\,\mathrm J},
\]
\[
\boxed{Q_p=n\frac72 R\Delta T\approx-1.985\times10^3\,\mathrm J},
\quad
\boxed{W_{\mathrm{on},p}\approx567\,\mathrm J}.
\]
\crosscheck 两过程终态不同：等温压缩末温保持273 K，等压压缩末温为136.5 K。
两过程均是压缩，所以气体对外功为负、外界对气体做功为正。

\sourceexample{8.2}{112}{120}
\problemtext $\nu$ mol刚性双原子理想气体从外界吸热$Q$并对外做功$W$，
求单个分子的平均\emph{平动}动能增加量，使用$N_A$表示阿伏伽德罗常量。
\workings
刚性双原子分子的总有效自由度$i=5$，由第一定律
\[
Q-W=\Delta U=\nu\frac52 R\Delta T
=\nu\frac52N_Ak_{\mathrm B}\Delta T.
\]
每分子的平均平动动能为$\bar\varepsilon_{\mathrm{tr}}=\frac32k_{\mathrm B}T$，
因此
\[
\boxed{\Delta\bar\varepsilon_{\mathrm{tr}}
=\frac32k_{\mathrm B}\Delta T
=\frac{3(Q-W)}{5\nu N_A}}.
\]
\crosscheck 题问\emph{单分子平动动能}，不是整团气体的总内能；
$Q,W$均用焦耳时，答案单位也是每分子焦耳。

\sourceexample{8.3}{112--113}{120--121}
\problemtext 1 mol刚性多原子分子理想气体，有效自由度$i=6$，
初始$1.0$ atm、$27^\circ$C，准静态绝热压缩至16 atm。
求内能增加、气体对外做功和末态分子数密度。
\workings
按原题指定$i=6$，定容热容$C_{V,m}=3R$，
$\gamma=(i+2)/i=4/3$。
可逆绝热温压关系为
\[
T_2=T_1\left(\frac{p_2}{p_1}\right)^{(\gamma-1)/\gamma}
\approx 300(16)^{1/4}=600\,\mathrm K.
\]
所以
\[
\boxed{\Delta U=3R(600-300)\approx7.48\times10^3\,\mathrm J},
\quad
\boxed{W_{\mathrm{gas}}=-\Delta U\approx-7.48\times10^3\,\mathrm J}.
\]
由$p=nk_{\mathrm B}T$得到末态数密度
\[
\boxed{n_2=\frac{16(1.013\times10^5)}{k_{\mathrm B}(600)}
\approx1.96\times10^{26}\,\mathrm{m^{-3}}}.
\]
\crosscheck $Q=0$，气体被压缩时自身对外功为负，而内能增加。
题面明确$i=6$，不应把“多原子”擅改为刚性双原子的$i=5$。

\sourceexample{8.4}{113}{121}
\problemtext 1 mol理想气体在$T_h=400\,\mathrm K$和$T_c=300\,\mathrm K$之间做可逆卡诺循环，
高温等温膨胀时体积由$V_1=0.001\,\mathrm{m^3}$变为
$V_2=0.005\,\mathrm{m^3}$，求每循环吸热、净功、向冷源放热。
\workings
高温等温膨胀时内能不变，吸热全部化为对外功：
\[
\boxed{Q_h=RT_h\ln\frac{V_2}{V_1}
=400R\ln5\approx5.35\times10^3\,\mathrm J}.
\]
卡诺效率只与热源温度有关，
\[
\eta=1-\frac{T_c}{T_h}=\frac14,\qquad
\boxed{W_{\mathrm{net}}=\eta Q_h\approx1.34\times10^3\,\mathrm J},
\]
\[
\boxed{Q_c=Q_h-W_{\mathrm{net}}\approx4.01\times10^3\,\mathrm J}.
\]
\crosscheck 每循环内能改变量为零，故$Q_h=Q_c+W_{\mathrm{net}}$；
$Q_c$这里是放热量的\emph{正大小}。

\sourceexample{8.5}{113--114}{121--122}
\problemtext 1 mol刚性双原子理想气体沿$p$--$V$图上三角形
$A(8.2\,\mathrm L,1\,\mathrm{atm})\to
B(16.4\,\mathrm L,3\,\mathrm{atm})\to
C(16.4\,\mathrm L,1\,\mathrm{atm})\to A$做循环。
求各状态温度、净功、$A\to B$的平均摩尔热容和热效率。
\workings
取$R\simeq8.31\,\mathrm{J/(mol\,K)}$及$1\,\mathrm{atm}\simeq1.013\times10^5\,\mathrm{Pa}$。
由$pV=RT$，
\[
\boxed{T_A\approx100\,\mathrm K},\quad
\boxed{T_B\approx600\,\mathrm K},\quad
\boxed{T_C\approx200\,\mathrm K}.
\]
循环净功等于三角形面积，
\[
\boxed{W_{\mathrm{net}}=\frac12(3-1)p_A(16.4-8.2)\,\mathrm L
=p_AV_A\approx831\,\mathrm J}.
\]
$A\to B$的内能增加和对外功分别是
\[
\Delta U_{AB}=\frac52R(T_B-T_A)=\frac{25}{2}RT_A,\quad
W_{AB}=\frac{p_A+p_B}{2}(V_B-V_A)=2RT_A.
\]
所以吸热$Q_{AB}=\frac{29}{2}RT_A$，
\[
\boxed{C_{AB}=\frac{Q_{AB}}{T_B-T_A}=\frac{29}{10}R}.
\]
其他两段均放热，整个循环只在$A\to B$吸热，故
\[
\boxed{\eta=\frac{W_{\mathrm{net}}}{Q_{AB}}
=\frac{2}{29}\approx6.90\%}.
\]
\crosscheck $B\to C$等容不做功，$C\to A$等压压缩做负功；
三角形沿顺时针方向运行，其净功为正。

\sourceexample{8.6}{114}{122}
\problemtext 绝热刚性容器被隔板分成体积$V_1,V_2$的两室，
分别装两种\emph{不同种类}的理想气体；初态两侧具有相同温度$T$和压强$p$。
取下隔板并使两气体无化学反应地均匀混合，求总熵变。
\workings
混合前两部分摩尔数分别为$\nu_i=pV_i/(RT)$。
忽略化学反应与相互作用混合热，最终温度仍为$T$，
每种气体在等温意义下由$V_i$“自由扩展”至$V_1+V_2$。
状态函数熵变可用假想可逆等温膨胀路径计算：
\[
\Delta S_i=\nu_iR\ln\frac{V_1+V_2}{V_i}.
\]
两者相加得
\[
\boxed{\Delta S=p\frac{V_1}{T}\ln\frac{V_1+V_2}{V_1}
+p\frac{V_2}{T}\ln\frac{V_1+V_2}{V_2}}>0.
\]
\crosscheck 此式针对可区分的\emph{不同种类}理想气体；
若是相同气体且初态$T,p$均相同，撤除隔板不产生经典混合熵，
不可忽略气体相同性而直接套用上式。
% === END guide-ch08-examples.tex ===
